\documentclass[../../main.tex]{subfiles}

\begin{document}

\section{一般的 Lebesgue 积分}\label{section:ch04-一般积分}

对于 $E$ 上的扩充实值函数 $f$, 我们分别定义 $f$ 的正部 $f^{+}$ 与负部 $f^{-}$ 为
\begin{align*}
f^{+}(x)=\max\{f(x),0\}\quad\text{与}\quad f^{-}(x)=\max\{-f(x),0\},\quad x\in E.
\end{align*}
则 $f^{+}$ 和 $f^{-}$ 是 $E$ 上的非负函数, 且
\begin{align*}
\text{在 }E\text{ 上},\quad f=f^{+}-f^{-},\quad |f|=f^{+}+f^{-}.
\end{align*}
观察到 $f$ 是可测的当且仅当 $f^{+}$ 和 $f^{-}$ 都是可测的.

\begin{proposition}\label{proposition:ch04-正负部可积}
令 $f$ 为 $E$ 上的可测函数. 则 $f^{+}$ 和 $f^{-}$ 在 $E$ 上可积当且仅当 $|f|$ 在 $E$ 上可积.
\end{proposition}

\begin{proof}
假设 $f^{+}$ 和 $f^{-}$ 是非负可积函数. 根据关于非负函数的积分的线性, $|f|=f^{+}+f^{-}$ 在 $E$ 上可积. 反之, 假定 $|f|$ 在 $E$ 上可积. 由于在 $E$ 上 $0\leqslant f^{+}\leqslant|f|$ 且 $0\leqslant f^{-}\leqslant|f|$, 我们从非负函数的积分的单调性推出 $f^{+}$ 和 $f^{-}$ 在 $E$ 上可积.
\end{proof}

\begin{definition}[Lebesgue 可积函数]\label{definition:ch04-可积}
在 $E$ 上的可测函数 $f$ 称为在 $E$ 上\textbf{是可积的}, 若 $|f|$ 在 $E$ 上可积. 当其成立时, 我们定义 $f$ 在 $E$ 上的积分为
\begin{align*}
\int_{E}f=\int_{E}f^{+}-\int_{E}f^{-}.
\end{align*}
\end{definition}

当然, 对于非负函数 $f$, 由于在 $E$ 上 $f=f^{+}$ 而 $f^{-}\equiv 0$, 积分的这个定义与我们刚考虑的那个定义一致. 根据有限支撑的有界可测函数的积分的线性, 以上积分的定义也与该函数类的积分的定义一致.

\begin{proposition}\label{proposition:ch04-可积a.e.有限}
令 $f$ 在 $E$ 上可积. 则 $f$ 在 $E$ 上 a.e. 有限且
\begin{align}
\text{若 }E_{0}\subseteq E\text{ 且 }m(E_{0})=0,\quad\int_{E}f=\int_{E\sim E_{0}}f.\label{align:ch04-可积零测}
\end{align}
\end{proposition}

\begin{proof}
命题\cref{proposition:ch04-可积有限}告诉我们 $|f|$ 在 $E$ 上 a.e. 有限. 因此 $f$ 在 $E$ 上 a.e. 有限. 此外, 将\eqref{align:ch04-零测忽略}应用于 $f$ 的正部与负部得到了\eqref{align:ch04-可积零测}.
\end{proof}

以下关于可积性的判别法是实数级数收敛的比较判别法在 Lebesgue 积分上的对应结果.

\begin{proposition}[积分的比较判别法]\label{proposition:ch04-比较判别}
令 $f$ 为 $E$ 上的可测函数. 假定存在非负函数 $g$, 它在 $E$ 上可积且在
\begin{align*}
E\text{ 上 }|f|\leqslant g
\end{align*}
的意义下控制 $f$. 则 $f$ 在 $E$ 上可积且
\begin{align*}
\left|\int_{E}f\right|\leqslant\int_{E}|f|.
\end{align*}
\end{proposition}

\begin{proof}
根据非负函数的积分的单调性, $|f|$ 可积, 且因此 $f$ 是可积的. 根据实数的三角不等式与非负函数的积分的线性,
\begin{align*}
\left|\int_{E}f\right|=\left|\int_{E}f^{+}-\int_{E}f^{-}\right|\leqslant\int_{E}f^{+}+\int_{E}f^{-}=\int_{E}|f|.
\end{align*}
\end{proof}

我们已到了研究一元实变函数的一般性的 Lebesgue 积分的最后步骤. 在证明积分的线性性质之前, 我们需要处理关于积分的一个难点, 该难点在讨论可测性时已处理过. 即对于在 $E$ 上可积的两个函数 $f$ 与 $g$, 它们的和 $f+g$ 在 $E$ 中使得 $f$ 与 $g$ 取反号的无穷值的点不是适当定义的. 然而, 根据命题\cref{proposition:ch04-可积a.e.有限}, 若我们定义 $A$ 为 $E$ 中使得 $f$ 和 $g$ 都有限的那些点的集合, 则 $m(E\sim A)=0$. 一旦我们证明了 $f+g$ 在 $A$ 上可积, 我们定义
\begin{align*}
\int_{E}(f+g)=\int_{A}(f+g).
\end{align*}
我们从\eqref{align:ch04-可积零测}推出 $\int_{E}(f+g)$ 等于任何 $(f+g)|_{A}$ 在 $E$ 上的延拓的扩充实值函数在 $E$ 上的积分.

\begin{theorem}[积分的线性与单调性]\label{theorem:ch04-一般线性}
令函数 $f$ 和 $g$ 在 $E$ 上可积. 则对任何 $\alpha$ 和 $\beta$, 函数 $\alpha f+\beta g$ 在 $E$ 上可积且
\begin{align*}
\int_{E}(\alpha f+\beta g)=\alpha\int_{E}f+\beta\int_{E}g.
\end{align*}
此外,
\begin{align*}
\text{若在 }E\text{ 上 }f\leqslant g,\quad\text{则}\int_{E}f\leqslant\int_{E}g.
\end{align*}
\end{theorem}

\begin{proof}
若 $\alpha>0$, 则 $[\alpha f]^{+}=\alpha f^{+}$ 且 $[\alpha f]^{-}=\alpha f^{-}$, 而若 $\alpha<0$, $[\alpha f]^{+}=-\alpha f^{-}$ 且 $[\alpha f]^{-}=-\alpha f^{+}$. 因此, $\int_{E}\alpha f=\alpha\int_{E}f$, 这是由于我们对非负函数 $f$ 和 $\alpha>0$ 证明了这个结论. 因此仅须在 $\alpha=\beta=1$ 的情形下证明线性. 根据非负函数的积分的线性, $|f|+|g|$ 在 $E$ 上可积. 由于在 $E$ 上 $|f+g|\leqslant|f|+|g|$, 根据积分的比较判别法, $f+g$ 在 $E$ 上也是可积的. 命题\cref{proposition:ch04-可积a.e.有限}告诉我们 $f$ 和 $g$ 在 $E$ 上是 a.e. 有限的. 根据同一命题, 通过从 $E$ 中去掉一个可能的测度为零的集合, 我们可以假设 $f$ 和 $g$ 在 $E$ 上有限. 要证明线性即要证明
\begin{align}
\int_{E}[f+g]^{+}-\int_{E}[f+g]^{-}=\left[\int_{E}f^{+}-\int_{E}f^{-}\right]+\left[\int_{E}g^{+}-\int_{E}g^{-}\right].\label{align:ch04-线性证明}
\end{align}
但在 $E$ 上,
\begin{align*}
(f+g)^{+}-(f+g)^{-}=f+g=(f^{+}-f^{-})+(g^{+}-g^{-}).
\end{align*}
因此, 由于这 $6$ 个函数的每个在 $E$ 上取实值,
\begin{align*}
\text{在 }E\text{ 上},\quad(f+g)^{+}+f^{-}+g^{-}=(f+g)^{-}+f^{+}+g^{+}.
\end{align*}
我们从非负函数的线性推出
\begin{align*}
\int_{E}(f+g)^{+}+\int_{E}f^{-}+\int_{E}g^{-}=\int_{E}(f+g)^{-}+\int_{E}f^{+}+\int_{E}g^{+}.
\end{align*}
由于 $f$, $g$ 和 $f+g$ 在 $E$ 上可积, 这 $6$ 个积分的每个是有限的. 重新排列这些积分以得到\eqref{align:ch04-线性证明}, 这完成了线性的证明.

为证明单调性我们可以假设 $g$ 和 $f$ 在 $E$ 上有限. 在 $E$ 上定义 $h=g-f$. 则 $h$ 是 $E$ 上的正确定义的非负可测函数. 根据可积函数的积分的线性和非负函数积分的单调性,
\begin{align*}
\int_{E}g-\int_{E}f=\int_{E}(g-f)=\int_{E}h\geqslant 0.
\end{align*}
\end{proof}

\begin{corollary}[积分在区域上的可加性]\label{corollary:ch04-一般可加}
令 $f$ 在 $E$ 上可积. 假设 $A$ 和 $B$ 是 $E$ 的不交可测子集. 则
\begin{align}
\int_{A\cup B}f=\int_{A}f+\int_{B}f.\label{align:ch04-一般可加}
\end{align}
\end{corollary}

\begin{proof}
观察到在 $E$ 上 $|f\cdot\chi_{A}|\leqslant|f|$ 且 $|f\cdot\chi_{B}|\leqslant|f|$. 根据积分的比较判别法, 可测函数 $f\cdot\chi_{A}$ 和 $f\cdot\chi_{B}$ 在 $E$ 上可积. 由于 $A$ 和 $B$ 不交
\begin{align}
\text{在 }E\text{ 上},\quad f\cdot\chi_{A\cup B}=f\cdot\chi_{A}+f\cdot\chi_{B}.\label{align:ch04-特征分解}
\end{align}
但对 $E$ 的任何可测子集 $C$ (见习题 28),
\begin{align*}
\int_{C}f=\int_{E}f\cdot\chi_{C}.
\end{align*}
因此\eqref{align:ch04-一般可加}从\eqref{align:ch04-特征分解}和积分的线性性质得出.
\end{proof}

有界收敛定理的以下推广提供了积分号下取极限的另一种合理性准则.

\begin{theorem}[Lebesgue 控制收敛定理]\label{theorem:ch04-LDCT}
令 $\{f_{n}\}$ 为 $E$ 上的可测函数序列. 假定存在 $E$ 上的可积函数 $g$ 在
\begin{align*}
\text{对所有 }n,\quad\text{在 }E\text{ 上 }|f_{n}|\leqslant g
\end{align*}
的意义下控制 $\{f_{n}\}$.
\begin{align*}
\text{若在 }E\text{ 上 a.e. 逐点 }\{f_{n}\}\to f,\text{ 则 }f\text{ 在 }E\text{ 上可积且}\lim_{n\to\infty}\int_{E}f_{n}=\int_{E}f.
\end{align*}
\end{theorem}

\begin{proof}
由于在 $E$ 上 $|f_{n}|\leqslant g$ 且在 $E$ 上 a.e. $|f|\leqslant g$, 而 $g$ 在 $E$ 上可积, 根据积分的比较判别法, $f$ 和每个 $f_{n}$ 在 $E$ 上也是可积的. 通过从 $E$ 去掉测度为零的集合的可数族以及利用 Lebesgue 测度的可数可加性, 我们从命题\cref{proposition:ch04-可积a.e.有限}得知, 可以假设 $f$ 和每个 $f_{n}$ 在 $E$ 上是有限的. 函数 $g-f$ 以及对每个 $n$ 函数 $g-f_{n}$ 是恰当定义的、非负可测的. 此外, 序列 $\{g-f_{n}\}$ 在 $E$ 上 a.e. 逐点收敛于 $g-f$. Fatou 引理告诉我们
\begin{align*}
\int_{E}(g-f)\leqslant\liminf\int_{E}(g-f_{n}).
\end{align*}
因此, 根据可积函数的积分的线性性质,
\begin{align*}
\int_{E}g-\int_{E}f=\int_{E}(g-f)\leqslant\liminf\int_{E}(g-f_{n})=\int_{E}g-\limsup\int_{E}f_{n},
\end{align*}
即
\begin{align*}
\limsup\int_{E}f_{n}\leqslant\int_{E}f.
\end{align*}
类似地, 考虑序列 $\{g+f_{n}\}$, 我们得到
\begin{align*}
\int_{E}f\leqslant\liminf\int_{E}f_{n}.
\end{align*}
证毕.
\end{proof}

Lebesgue 控制收敛定理的以下推广常常是有用的 (见习题 33). 它的证明我们留作练习 (见习题 32).

\begin{theorem}[一般的 Lebesgue 控制收敛定理]\label{theorem:ch04-一般LDCT}
令 $\{f_{n}\}$ 为 $E$ 上的 a.e. 逐点收敛于 $f$ 的 $E$ 上的可测函数序列. 假定存在 $E$ 上的非负可测函数序列 $\{g_{n}\}$, 它在 $E$ 上 a.e. 逐点收敛于 $g$ 且在
\begin{align*}
\text{对所有 }n,\quad\text{在 }E\text{ 上 }|f_{n}|\leqslant g_{n}
\end{align*}
的意义下控制 $\{f_{n}\}$.
\begin{align*}
\text{若}\lim_{n\to\infty}\int_{E}g_{n}=\int_{E}g<\infty,\text{则}\lim_{n\to\infty}\int_{E}f_{n}=\int_{E}f.
\end{align*}
\end{theorem}

\begin{remark}
在 Fatou 引理和 Lebesgue 控制收敛定理中, 在 $E$ 上 a.e. 逐点收敛而非在整个 $E$ 上收敛的假设不是一般性的修饰. 它对于这些结果将来的应用是必要的. 我们提供这个必要性的说明. 假定 $f$ 是整个 $\mathbf{R}$ 上的递增函数. 即将介绍的一个 Lebesgue 定理 (第 6 章的 Lebesgue 定理) 告诉我们
\begin{align}
\text{对几乎所有 }x,\quad\lim_{n\to\infty}\frac{f(x+1/n)-f(x)}{1/n}=f'(x)<\infty.\label{align:ch04-导数极限}
\end{align}
从这与 Fatou 引理我们将证明, 对任何闭有界区间 $[a,b]$,
\begin{align*}
\int_{a}^{b}f'\leqslant f(b)-f(a).
\end{align*}
一般地, 给定非退化闭有界区间 $[a,b]$ 和 $[a,b]$ 的测度为零的子集 $A$, 存在 $[a,b]$ 上的递增函数 $f$ 使得在 $A$ 中的每一点, \eqref{align:ch04-导数极限}中的极限不存在 (见第 6 章的习题 10).
\end{remark}

\begin{exercise}\label{exercise:ch04-24}
令 $f$ 为 $E$ 上的非负可测函数.
\begin{enumerate}[(i)]
\item 证明存在 $E$ 上的非负简单函数的递增序列 $\{\varphi_{n}\}$, 该序列的每个函数有有限支撑, 它在 $E$ 上逐点收敛于 $f$.
\item 证明 $\int_{E}f=\sup\{\int_{E}\varphi\mid\varphi\text{ 是简单的, 具有有限支撑, 且在 }E\text{ 上 }0\leqslant\varphi\leqslant f\}$.
\end{enumerate}
\end{exercise}

\begin{exercise}\label{exercise:ch04-25}
令 $\{f_{n}\}$ 为 $E$ 上的非负可测函数序列, 该序列在 $E$ 上逐点收敛于 $f$. 假定在 $E$ 上对每个 $n$, $f_{n}\leqslant f$. 证明
\begin{align*}
\lim_{n\to\infty}\int_{E}f_{n}=\int_{E}f.
\end{align*}
\end{exercise}

\begin{exercise}\label{exercise:ch04-26}
证明单调收敛定理对递减的函数序列可能不成立.
\end{exercise}

\begin{exercise}\label{exercise:ch04-27}
证明 Fatou 引理的以下推广: 若 $\{f_{n}\}$ 是 $E$ 上的非负可测函数序列, 则
\begin{align*}
\int_{E}\liminf f_{n}\leqslant\liminf\int_{E}f_{n}.
\end{align*}
\end{exercise}

\begin{exercise}\label{exercise:ch04-28}
令 $f$ 在 $E$ 上可积而 $C$ 是 $E$ 的可测子集. 证明 $\int_{C}f=\int_{E}f\cdot\chi_{C}$.
\end{exercise}

\begin{exercise}\label{exercise:ch04-29}
对于 $[1,\infty)$ 上的在有界集上有界的可测函数 $f$, 对每个自然数 $n$, 定义 $a_{n}=\int_{n}^{n+1}f$. $f$ 在 $[1,\infty)$ 上可积当且仅当级数 $\sum_{n=1}^{\infty}a_{n}$ 收敛, 该结论是否成立? $f$ 在 $[1,\infty)$ 上可积当且仅当级数 $\sum_{n=1}^{\infty}a_{n}$ 绝对收敛, 该结论是否成立?
\end{exercise}

\begin{exercise}\label{exercise:ch04-30}
令 $g$ 为 $E$ 上的非负可积函数, 且假定 $\{f_{n}\}$ 是 $E$ 上的可测函数序列, 使得对每个 $n$, 在 $E$ 上 a.e. $|f_{n}|\leqslant g$. 证明
\begin{align*}
\int_{E}\liminf f_{n}\leqslant\liminf\int_{E}f_{n}\leqslant\limsup\int_{E}f_{n}\leqslant\int_{E}\limsup f_{n}.
\end{align*}
\end{exercise}

\begin{exercise}\label{exercise:ch04-31}
令 $f$ 为 $E$ 上的可测函数, 它在 $E$ 上可表示为 $f=g+h$, 其中 $g$ 在 $E$ 上有限且可积而 $h$ 在 $E$ 上非负. 定义 $\int_{E}f=\int_{E}g+\int_{E}h$. 证明这基于以下意义是正确定义的: 它独立于和为 $f$ 的有限可积函数 $g$ 和非负函数 $h$ 的特定选取.
\end{exercise}

\begin{exercise}\label{exercise:ch04-32}
遵循 Lebesgue 控制收敛定理的证明方法, 证明一般的 Lebesgue 控制收敛定理, 用序列 $\{g_{n}-f_{n}\}$ 和 $\{g_{n}+f_{n}\}$ 分别代替 $\{g-f_{n}\}$ 和 $\{g+f_{n}\}$.
\end{exercise}

\begin{exercise}\label{exercise:ch04-33}
令 $\{f_{n}\}$ 为 $E$ 上的可积函数序列, 使得在 $E$ 上 a.e. $f_{n}\to f$ 而 $f$ 在 $E$ 上可积. 证明 $\int_{E}|f-f_{n}|\to 0$ 当且仅当 $\lim_{n\to\infty}\int_{E}|f_{n}|=\int_{E}|f|$. (提示: 用一般的 Lebesgue 控制收敛定理.)
\end{exercise}

\begin{exercise}\label{exercise:ch04-34}
令 $f$ 为 $\mathbf{R}$ 上的非负可测函数. 证明
\begin{align*}
\lim_{n\to\infty}\int_{-n}^{n}f=\int_{\mathbf{R}}f.
\end{align*}
\end{exercise}

\begin{exercise}\label{exercise:ch04-35}
令 $f$ 为定义在正方形 $Q=\{(x,y)\mid 0\leqslant x\leqslant 1,\ 0\leqslant y\leqslant 1\}$ 上的双变量 $(x,y)$ 实值函数, 且对每个固定的 $y$ 值它是 $x$ 的可测函数. 假定对每个固定的 $x$ 值, $\lim_{y\to 0}f(x,y)=f(x)$, 且对所有 $y$, 我们有 $|f(x,y)|\leqslant g(x)$, 其中 $g$ 在 $[0,1]$ 上可积. 证明
\begin{align*}
\lim_{y\to 0}\int_{0}^{1}f(x,y)\,dx=\int_{0}^{1}f(x)\,dx.
\end{align*}
也证明: 若函数 $f(x,y)$ 对每个 $x$ 在 $y$ 连续, 则
\begin{align*}
h(y)=\int_{0}^{1}f(x,y)\,dx
\end{align*}
是 $y$ 的连续函数.
\end{exercise}

\begin{exercise}\label{exercise:ch04-36}
令 $f$ 为定义在正方形 $Q=\{(x,y)\mid 0\leqslant x\leqslant 1,\ 0\leqslant y\leqslant 1\}$ 上的双变量 $(x,y)$ 实值函数, 且对每个固定的 $y$ 值它是 $x$ 的可测函数. 对每个 $(x,y)\in Q$, 令偏导数 $\partial f/\partial y$ 存在. 假定存在函数 $g$ 在 $[0,1]$ 上可积, 使得对所有 $(x,y)\in Q$, $|\partial f/\partial y(x,y)|\leqslant g(x)$. 证明
\begin{align*}
\text{对所有 }y\in[0,1],\quad\frac{d}{dy}\left[\int_{0}^{1}f(x,y)\,dx\right]=\int_{0}^{1}\frac{\partial f}{\partial y}(x,y)\,dx.
\end{align*}
\end{exercise}

\end{document}